%! AP Statistics — Unit 3, Lesson 3.5 — Student Guided Notes
\documentclass[10pt]{article}
\usepackage{apstats-article}
\usepackage{apstats-boxes}

\begin{document}

\pageheader{Unit 3, Lesson 3.5}{Guided Notes}
\namedateperiod

\vspace{0.12in}

\begin{objectivebox}
\small By the end of this lesson, I will be able to\dots\\[4pt]
\writeline\\[1pt]
\writeline\\[1pt]
\writeline
\end{objectivebox}

\vspace{0.1in}

\begin{vocabbox}
\termblanklong{Experimental unit}
\termblanklong{Factor (explanatory variable in an experiment)}
\termblanklong{Treatment}
\termblanklong{Response variable}
\termblanklong{Control group / Placebo}
\termblanklong{Random assignment}
\termblanklong{Blocking variable}
\termblanklong{Matched pairs design}
\end{vocabbox}

\vspace{0.1in}

\begin{notesbox}{Components of an Experiment (VAR-3.A)}
\small
\renewcommand{\arraystretch}{1.8}
\begin{tabularx}{\linewidth}{@{} >{\bfseries\color{navy}}p{0.3\linewidth} X @{}}
Experimental units & The \blank{6cm} assigned to treatments. \\
Factor & The \blank{3cm} variable whose levels are \blank{3cm} intentionally. \\
Treatment & A specific \blank{3cm} (or combination of levels) of the factor. \\
Response variable & The \blank{3cm} measured after treatments are applied. \\
Confounding variable & Related to both the \blank{3cm} and the \blank{3cm};
may create a false association. \\
\end{tabularx}
\end{notesbox}

\vspace{0.1in}

\begin{notesbox}{Principles of a Well-Designed Experiment (VAR-3.B.1)}
\small
A well-designed experiment must include:
\begin{enumerate}[leftmargin=1.5em, itemsep=6pt]
  \item \textbf{Comparison:} at least \blank{2cm} treatment groups (one may be a \blank{3cm}).
  \item \textbf{Random assignment:} units assigned to treatments \blank{4cm}; this \blank{3cm}
        confounding variables and allows \blank{3cm} conclusions.
  \item \textbf{Replication:} more than \blank{2cm} unit per treatment group.
  \item \textbf{Control:} potential \blank{4cm} variables are managed (by design or by
        holding constant).
\end{enumerate}
\end{notesbox}

\vspace{0.1in}

\begin{notesbox}{Experimental Designs (VAR-3.C)}
\small
\renewcommand{\arraystretch}{1.9}
\begin{tabularx}{\linewidth}{@{} >{\bfseries\color{navy}}p{0.26\linewidth}
    >{\raggedright\arraybackslash}X @{}}
Completely randomized (CRD) &
All units randomly assigned to treatments --- \blank{4cm}. \\
Block design &
Units grouped into \blank{4cm} based on a blocking variable; treatments assigned
\blank{3cm} within each block. Reduces \blank{3cm} due to the blocking variable. \\
Matched pairs &
A special block design; blocks of \blank{2cm} matched on a key variable (or same unit
gets \blank{3cm} treatments in random order). \\
Single-blind &
\blank{4cm} don't know which treatment they received. \\
Double-blind &
Neither \blank{4cm} nor \blank{5cm} know treatment assignments. \\
\end{tabularx}
\end{notesbox}

\vspace{0.1in}

\begin{practicebox}
\small A study tests a new energy drink on reaction time. 60 college students are
randomly assigned: 30 drink the energy drink, 30 drink a look-alike water.
Neither students nor the researcher measuring reaction time know which was which.
\begin{enumerate}[leftmargin=1.5em, itemsep=6pt]
  \item Experimental units: \blank{8cm}
  \item Factor: \blank{8cm}
  \item Treatments: \blank{8cm}
  \item Response variable: \blank{8cm}
  \item Type of blinding: \blank{8cm}
  \item Design type: \blank{8cm}
\end{enumerate}
\end{practicebox}

\end{document}
